%%=============================================================================
%% Voorwoord
%%=============================================================================

\chapter*{\IfLanguageName{dutch}{Woord vooraf}{Preface}}%
\label{ch:voorwoord}

%% TODO:
%% Het voorwoord is het enige deel van de bachelorproef waar je vanuit je
%% eigen standpunt (``ik-vorm'') mag schrijven. Je kan hier bv. motiveren
%% waarom jij het onderwerp wil bespreken.
%% Vergeet ook niet te bedanken wie je geholpen/gesteund/... heeft

Tijdens mijn studie heb ik zelf ervaren hoe uitdagend het is om met concentratieproblemen om te gaan. Taken plannen, prioriteiten stellen en gefocust blijven zijn voor velen vanzelfsprekend, maar voor mij en voor veel anderen vraagt dit vaak extra inspanning. Deze persoonlijke ervaring heeft mij gemotiveerd om dit onderwerp verder te onderzoeken en er mijn bachelorproef aan te wijden.

Met deze bachelorproef wil ik niet alleen mijn eigen inzichten verdiepen, maar ook anderen helpen die met gelijkaardige moeilijkheden kampen. Daarom heb ik gewerkt aan het ontwikkelen van een AI-gebaseerd taakallocatiemodel. Dit model heeft als doel taken beter af te stemmen op de sterktes en voorkeuren van een werknemer, zodat de werkdruk beter beheersbaar wordt en stress vermindert. Door werknemers taken te laten uitvoeren die hen beter liggen, hoop ik bij te dragen aan een meer haalbare en aangenamere manier van werken in een organisatie.

Tot slot wil ik graag enkele mensen bedanken die mij gesteund hebben tijdens dit proces: